\documentclass{llncs}
\usepackage{graphicx}
\usepackage{amsmath,amssymb}
\usepackage{booktabs}
\usepackage{hyperref}

\begin{document}

\title{Paper Title for Springer LNCS}
\titlerunning{Abbreviated Paper Title}

\author{First Author\inst{1}\orcidID{0000-1111-2222-3333} \and
Second Author\inst{2,3}\orcidID{1111-2222-3333-4444}}

\authorrunning{F. Author et al.}

\institute{Princeton University, Princeton NJ 08544, USA \and
Springer Heidelberg, Tiergartenstr. 17, 69121 Heidelberg, Germany
\email{lncs@springer.com}\\
\url{http://www.springer.com/gp/computer-science/lncs} \and
ABC Institute, Rupert-Karls-University Heidelberg, Heidelberg, Germany\\
\email{\{abc,lncs\}@uni-heidelberg.de}}

\maketitle

\begin{abstract}
The abstract should summarize the contents of the paper within 150--250 words.

\keywords{First keyword \and Second keyword \and Another keyword.}
\end{abstract}

\section{Introduction}
Springer LNCS format requires standard sectioning and author affiliations, referencing relevant work~\cite{sample}.

\section{Methodology}
Mathematical descriptions and proofs.
\begin{equation}
\mathcal{L}_{total} = \mathcal{L}_{task} + \lambda \mathcal{L}_{reg}
\end{equation}

\section{Results}
\begin{table}[ht]
\centering
\caption{Benchmark evaluations}\label{tab1}
\begin{tabular}{lll}
\hline
Heading level & Example & Font size and style\\
\hline
Title (centered) & {\Large\bfseries Lecture Notes} & 14 point, bold\\
1st-level heading & {\large\bfseries 1 Introduction} & 12 point, bold\\
\hline
\end{tabular}
\end{table}

\section{Conclusion}
Concluding remarks.

\bibliographystyle{splncs04}
\bibliography{references}

\end{document}
